\section{Method}
\label{sec:method}

\subsection{Preliminaries: GRPO Under Binary Rewards}
\label{sec:prelim}

Given a prompt $x$ and a policy $\pitheta$, GRPO \citep{shao2024deepseekmath} samples a group of $G$ completions $\{y_1, \ldots, y_G\} \sim \pitheta(\cdot|x)$ and computes binary verifier rewards $r_i = R(x, y_i) \in \{0, 1\}$.
The standard group-relative advantage is:
\begin{equation}
\label{eq:grpo_adv}
A_i^{\mathrm{GRPO}} = \frac{r_i - \bar{r}}{\sigma + \epsilon}, \quad \bar{r} = \frac{1}{G}\sum_{j=1}^G r_j, \quad \sigma = \mathrm{std}(\{r_j\}).
\end{equation}
The policy is updated via clipped surrogate loss with KL regularization against a reference policy $\piref$:
\begin{equation}
\Lcal_{\mathrm{PG}} = -\E_{x,y}\Big[\sum_t A_i \log \pitheta(y_t | x, y_{<t}) - \beta \, \mathrm{KL}(\pitheta \| \piref)\Big].
\end{equation}

\paragraph{The degeneracy problem.}
Under binary rewards with small $G$, a large fraction of groups are \emph{homogeneous}: all responses are correct ($r_i = 1$ for all $i$) or all are incorrect ($r_i = 0$ for all $i$).
For these groups, $\bar{r} \in \{0, 1\}$ and $\sigma = 0$, so $A_i^{\mathrm{GRPO}} = 0$ for every response in the group---producing zero gradient.
At $G{=}4$, if a prompt has success probability $p$, the probability of a homogeneous group is $p^4 + (1-p)^4$.
For a typical prompt with $p=0.3$, this gives $p^4 + (1-p)^4 = 0.0081 + 0.2401 = 0.248$, meaning about one quarter of groups produce no learning signal.
Dr.GRPO \citep{liu2025understanding} addresses length bias but retains this mean-centered formulation, so the degeneracy persists.

\subsection{Threshold-Anchored Signed Advantage (TASA)}
\label{sec:tasa}

We replace the group mean with a fixed \emph{threshold} $c$ as the advantage anchor.
For binary rewards, we set $c = 0.5$ (the midpoint between reward values 0 and 1).
The TASA advantage for the $i$-th response is:
\begin{equation}
\label{eq:tasa}
A_i^{\mathrm{TASA}} = \frac{(r_i - c)_+}{\Zp} - \frac{(c - r_i)_+}{\Zm},
\end{equation}
where $(x)_+ = \max(x, 0)$, $\Zp = \sum_{j=1}^G (r_j - c)_+$ is the total positive excess, and $\Zm = \sum_{j=1}^G (c - r_j)_+$ is the total negative deficit, both clamped to be at least $\epsilon > 0$.

\paragraph{Properties.}
TASA satisfies four desirable properties under binary rewards:
\begin{enumerate}[leftmargin=*,topsep=2pt,itemsep=1pt]
\item \textbf{Sign preservation.} If $r_i > c$, then $A_i > 0$; if $r_i \leq c$, then $A_i \leq 0$. Correct responses always receive positive advantage.
\item \textbf{Monotonicity.} $A_i$ is monotonically non-decreasing in $r_i$.
\item \textbf{Boundedness.} $A_i \in [-1, 1]$.
\item \textbf{Non-degeneracy.} Even all-correct groups ($r_i = 1$ for all $i$) produce $A_i = 1/G > 0$, and all-incorrect groups produce $A_i = -1/G < 0$. No group yields zero advantage.
\end{enumerate}

The non-degeneracy property is the key advantage over group-mean centering: TASA provides gradient signal for \emph{every} group, including the homogeneous ones that GRPO/Dr.GRPO discard.

\paragraph{Relationship to simpler baselines.}
Under binary rewards, TASA reduces to a count-normalized signed assignment: correct responses get $A_i = +1/n_+$ and incorrect ones get $A_i = -1/n_-$, where $n_+$ and $n_-$ are the counts of correct and incorrect responses.
A simpler alternative, $A_i = 2r_i - 1 \in \{-1, +1\}$, also preserves sign but ignores group composition.
TASA's normalization by group counts prevents the magnitude from growing when few responses share a label, which we find important for stable training (see \Cref{sec:exp}).

\paragraph{Potential failure modes.}
Because TASA always provides non-zero gradients, it can reinforce low-quality but technically correct solutions or penalize near-miss attempts that contain useful reasoning.
It also introduces a mild bias: in all-correct groups, every response is pushed up equally regardless of solution quality.
We view this as an acceptable tradeoff in the binary-reward setting, where the verifier provides no quality signal beyond correctness.

\subsection{Verified CE Replay}
\label{sec:replay}

To further exploit rare correct responses, we maintain a \emph{verified success bank} $\Bcal$ that stores successful completions encountered during training.
Each entry records the prompt $x$, completion token IDs $y^+$, and insertion step.
The bank is:
\begin{itemize}[leftmargin=*,topsep=2pt,itemsep=1pt]
\item \textbf{Per-prompt}: each prompt can hold at most $K$ distinct solutions (we use $K{=}2$).
\item \textbf{Hash-deduplicated}: new entries are rejected if their token-ID hash matches an existing entry for the same prompt.
\item \textbf{Recency-biased}: when the bank for a prompt is full, the oldest entry is replaced.
\end{itemize}

At each training step (after a warmup of $W$ steps), we sample a mini-batch of verified successes from $\Bcal$ and compute a supervised cross-entropy loss:
\begin{equation}
\label{eq:replay}
\Lcal_{\mathrm{replay}} = -\E_{(x, y^+) \sim \Bcal} \sum_t \log \pitheta(y^+_t | x, y^+_{<t}),
\end{equation}
where prompt tokens are masked (label $= {-}100$) so that only completion tokens contribute to the loss.

\paragraph{Total objective.}
The training loss combines TASA policy gradient, CE replay, and KL regularization:
\begin{equation}
\label{eq:total}
\Lcal = \Lcal_{\mathrm{TASA\text{-}PG}} + \lambda \, \Lcal_{\mathrm{replay}} + \beta \, \Lcal_{\mathrm{KL}},
\end{equation}
where $\lambda$ controls replay strength and $\beta$ controls KL penalty.

\paragraph{Comparison with PG-based replay.}
Both RePO \citep{li2025repo} and RLEP \citep{zhang2025rlep} replay past completions using the same policy-gradient objective with importance-sampling (IS) correction:
\[
\Lcal_{\mathrm{PG\text{-}replay}} = -\E_{(x,y)\sim\Bcal}\bigg[\sum_t \min\!\Big(\frac{\pitheta(y_t|x,y_{<t})}{\pi_{\mathrm{old}}(y_t|x,y_{<t})} A_i,\;\mathrm{clip}(\cdot)\Big)\bigg].
\]
This requires storing the generating policy's log-probabilities $\log\pi_{\mathrm{old}}$ for every replayed token and computing per-token IS ratios that become stale as the policy evolves.
Our CE replay (\Cref{eq:replay}) avoids both requirements: it stores only token IDs and uses a supervised objective that is agnostic to the generating policy.
\Cref{tab:replay_compare} summarizes the differences.

\begin{table}[t]
\centering
\caption{Comparison of replay loss designs. Our CE replay is simpler, requiring no stored log-probabilities or importance-sampling correction.}
\label{tab:replay_compare}
\small
\begin{tabular}{lccc}
\toprule
& \textbf{Our CE Replay} & \textbf{RePO} & \textbf{RLEP} \\
\midrule
Replay loss & Supervised CE & PG + IS clip & PG + IS clip \\
Stored data & Token IDs only & Token IDs + log-probs & Token IDs + log-probs \\
IS correction & None needed & Token-level IS ratio & Token-level IS ratio \\
Buffer scope & Per-prompt ($K{=}2$) & Global & Global \\
Deduplication & Hash-based & None & None \\
Online / Offline & Online & Online & Two-phase \\
\bottomrule
\end{tabular}
\end{table}
